% ============================================================
% ABSTRACT
% ============================================================
\chapter*{ABSTRACT}
\addcontentsline{toc}{chapter}{ABSTRACT}
\markboth{ABSTRACT}{ABSTRACT}

Dysarthria is a motor speech disorder that reduces the intelligibility of
spoken language by degrading the articulation of individual speech sounds.
Speech-and-language therapy for dysarthria depends on repeated, high-volume
articulation practice, but a clinician can only be present for a small fraction
of that practice. Existing automatic tools are of limited help: mainstream
automatic speech recognition returns a single word-level verdict, which tells a
learner \emph{that} a word was unclear but not \emph{which sound} failed, and
conventional mispronunciation-detection systems are trained on non-native
learner speech whose error patterns differ substantially from those of motor
speech disorders.

This dissertation presents \systemName, a phoneme-level pronunciation error
detection system for dysarthric and unclear speech, together with the deployable
web application built around it. The central methodological contribution is the
labelling strategy. A first attempt (Method~1) derived error labels by taking a
fixed percentile of the deviation distribution computed over \emph{all}
phonemes. Because that percentile was applied globally, a fixed fraction of
healthy control phonemes was forced into the error class by construction, and
the resulting model flagged errors in clear speech at inference time. Method~2,
the approach adopted in this work, redefines the ground truth around a
\emph{control baseline}: every phoneme produced by a healthy control speaker is
labelled correct by definition and retained as a positive anchor, while each
dysarthric phoneme is compared against its positionally matched control
realisation of the same target word, and is labelled as an error only when its
log-mel deviation falls in the top 30\% of the \emph{dysarthric} deviation
distribution. The error threshold is therefore never computed over control data,
and well-produced dysarthric phonemes are correctly retained as positive
examples.

The pipeline pairs 10{,}000 dysarthric--control recordings from the UASpeech
corpus. Both recordings of each pair are force-aligned against the same target
phoneme sequence, obtained from the control transcript via
grapheme-to-phoneme conversion, using a \texttt{wav2vec2} CTC phone recogniser
and the Viterbi forced-alignment algorithm. Each aligned phoneme is cropped with
20\,ms of context and converted into a fixed $64 \times 32$ log-mel patch,
yielding 98{,}306 labelled phoneme instances (14{,}746 error, 83{,}560 correct).
A CNN--BiLSTM classifier --- two convolutional blocks feeding a
128-unit bidirectional LSTM, 415{,}810 parameters in total --- is trained on a
class-balanced training split while the test split retains the natural class
distribution.

Five architectures were compared under a common protocol. CNN--BiLSTM achieved
the highest error-class recall (76.64\%) and the highest error-class
$F_1$ (74.04\%) at 73.13\% accuracy, and was selected because recall on the
error class is the operationally important quantity for a screening tool: a
missed error is a missed opportunity for corrective feedback. CNN--BiGRU
attained marginally higher accuracy (73.40\%) and precision (74.19\%) but lower
recall (71.75\%); the plain CNN baseline without a recurrent stage lost roughly
four points of $F_1$, confirming that the temporal modelling stage contributes
measurably. On the naturally imbalanced held-out split the selected model
reaches 71.99\% accuracy with 75\% recall and 32\% precision on the error class;
the low precision under natural imbalance is analysed in detail and attributed
jointly to the class prior and to the intrinsically noisy, distance-derived
nature of the labels. Per-phoneme analysis shows the highest predicted error
rates on the alveolar flap \ph{ɾ}, the rhotic schwa \ph{ɚ}, the velar nasal
\ph{ŋ} and reduced vowels, which is consistent with the articulatory profile of
dysarthria reported in the clinical literature.

The trained model was deployed as \systemName, a three-tier application
comprising a React progressive web application, a FastAPI inference service and
the model artefacts. The service reproduces the training-time feature pipeline
exactly, including the global mel normalisation constants, and adds two
production safeguards absent from the notebook: an audio quality gate that
rejects clipped, silent or noise-dominated recordings before inference, and a
Whisper-based target-word verification step that distinguishes ``a different
word was spoken'' from ``the target word was mispronounced''. Median end-to-end
latency for a single-word analysis is a few seconds on commodity CPU hardware,
and the whole system deploys as a single Docker image.

The work demonstrates that anchoring the ground truth to a healthy reference
speaker, rather than to a global statistic, is what turns a phoneme
error-detection model from a system that flags errors everywhere into one that
separates impaired from clear speech. Limitations --- most notably that the
labels are proxy labels derived from acoustic distance rather than clinician
annotations, and that the validation split was reused as the test split ---
are stated explicitly, and a programme of clinical validation is proposed as
future work.

\vspace{0.6cm}
\noindent\textbf{Keywords:} dysarthric speech; phoneme-level error detection;
mispronunciation detection and diagnosis; forced alignment; wav2vec2;
CNN--BiLSTM; log-mel spectrogram; UASpeech; computer-aided pronunciation
training; speech-language pathology.

\clearpage
